% ---------------------------------------------------------------------------
% DLPNO-IP-ADC(3): manuscript draft (theory and implementation only).
%
% Written for REVTeX 4-2 (AIP/JCP). If revtex4-2.cls is not installed the
% preamble falls back to `article' so that the file still compiles; the
% typeset result is then unbranded but complete.
% ---------------------------------------------------------------------------
\IfFileExists{revtex4-2.cls}{%
  \documentclass[aip,jcp,amsmath,amssymb,preprint,superscriptaddress]{revtex4-2}%
  \newcommand{\FALLBACK}{0}%
}{%
  \documentclass[11pt]{article}%
  \newcommand{\FALLBACK}{1}%
}
\if 1\FALLBACK
  \usepackage[margin=1in]{geometry}
  \usepackage{amsmath,amssymb,graphicx}
  % REVTeX-only markup, harmlessly absorbed so the file still typesets.
  \newcommand{\affiliation}[1]{}
  \newcommand{\email}[1]{}
\fi

\usepackage{bm}
\usepackage{booktabs}

% ------------------------------------------------------------- shorthands --
\newcommand{\ket}[1]{\left|#1\right\rangle}
\newcommand{\bra}[1]{\left\langle#1\right|}
\newcommand{\tps}{\tilde}
\newcommand{\pao}[1]{\tilde{#1}}
\newcommand{\Mss}{\mathbf{M}^{\text{ss}}}
\newcommand{\Msd}{\mathbf{M}^{\text{sd}}}
\newcommand{\Mds}{\mathbf{M}^{\text{ds}}}
\newcommand{\Mdd}{\mathbf{M}^{\text{dd}}}
\newcommand{\TCut}[1]{\textsc{TCut#1}}

\begin{document}

\title{A domain-based local pair natural orbital formulation of third-order
       algebraic diagrammatic construction theory for ionization potentials}

\author{}
\affiliation{Department of Chemistry, Indian Institute of Technology Bombay,
             Powai, Mumbai 400076, India}

\date{\today}

\begin{abstract}
We present a domain-based local pair natural orbital (DLPNO) formulation of
the algebraic diagrammatic construction scheme for the one-particle Green's
function through third order, DLPNO-IP-ADC(3), for closed-shell reference
determinants. The 2h1p space is spanned by pair natural orbitals constructed
over projected-atomic-orbital domains of localized occupied orbitals, so that
each occupied pair carries its own compact virtual space and no quantity is
ever expressed in the canonical virtual basis.

The central difficulty of a local ADC(3), which does not arise at second
order, is that the third-order self-energy and the first-order 2h1p--2h1p
block contain contractions in which a virtual index is summed between a
two-electron integral and an amplitude belonging to a \emph{different} occupied
pair. We show that such an index must be carried in a pair space that owns it
--- one containing the occupied label the virtual belongs to --- and that the
otherwise natural choice of projecting the amplitude into the target pair
introduces an error that does not vanish as the pair natural orbital threshold
is tightened. Every term of the ADC(3) matrix is classified accordingly, and
the routing that follows is given explicitly.

We further describe the three-step pair prescreening used to select strong
pairs, the population-based criterion used to build local fitting domains, the
metric that renders the spin-adapted secular problem symmetric, and a
composite correction that measures the truncation error of the doubles space
by folding a second-order, scaled-opposite-spin variant of the same equations
both canonically and locally. Two independent implementations are described:
a production code in C++ and an implementation of the same equations in dense
linear algebra, retained as a check on the working equations rather than on
their programming.
\end{abstract}

\maketitle

% ===========================================================================
\section{Introduction}
% ===========================================================================

The algebraic diagrammatic construction (ADC) scheme for the one-particle
Green's function provides a systematic, size-consistent and Hermitian
hierarchy of approximations to vertical ionization
potentials.\cite{Schirmer1982,Schirmer1983,Trofimov2005,Dreuw2015} Its
non-Dyson formulation replaces the frequency-dependent self-energy by a static
Hermitian secular matrix, so that ionization energies are obtained as
eigenvalues rather than as roots of a nonlinear equation, and a Davidson
procedure suffices. At second order, IP-ADC(2), the method is comparable in
cost to second-order M{\o}ller--Plesset theory and in accuracy to it; at third
order, IP-ADC(3), it approaches the accuracy of equation-of-motion coupled
cluster theory for the principal ionization
peaks,\cite{Trofimov2005,Dreuw2015} at an $O(N^6)$ cost that confines it to
small molecules.

The cost of any such method is dominated by the delocalization of the
canonical molecular orbitals. Local correlation
methods\cite{Pulay1983,Saebo1993} exploit the fact that the correlation energy
is a sum of essentially local contributions once the occupied orbitals are
localized: the correlation of a given occupied pair involves only virtual
functions in its neighbourhood, and pairs of distant occupied orbitals
contribute negligibly. Pair natural orbitals (PNOs) push this further by
compressing the virtual space of each pair to the eigenvectors of its own
first-order density,\cite{Neese2009} and the domain-based local pair natural
orbital (DLPNO) family combines the two: PNOs are constructed within
projected-atomic-orbital (PAO) domains chosen by orbital overlap, giving
methods whose cost grows linearly with system size while reproducing the
canonical result to within a controlled and systematically improvable
error.\cite{Riplinger2013a,Riplinger2013b,Pinski2015,Riplinger2016}

DLPNO ideas have been carried into excited- and ionized-state theory,
including the equation-of-motion coupled cluster treatment of ionization
potentials,\cite{Dutta2018} where the target quantity is an eigenvalue of an
effective Hamiltonian rather than an energy expectation value. The present
work applies them to IP-ADC. The second-order case is close to the
DLPNO-MP2 situation: the static self-energy is a contraction of first-order
amplitudes with exchange integrals, and both belong to the same pair.
Third order is qualitatively different. The third-order static self-energy
$\mathbf{M}^{\text{ss},(3)}$ contains fourteen distinct contractions, the coupling blocks
acquire third-order contributions, and the 2h1p--2h1p block acquires a
first-order coupling between different pairs. Several of these terms sum a
virtual index between a two-electron integral and an amplitude that belongs to
a different occupied pair, and it is precisely there that a naive
localization fails.

The failure is worth stating plainly, because it is silent. Every virtual
index in these expressions belongs to a definite occupied label: in
$t^{ij}_{ab}$ the index $a$ belongs to $i$ and $b$ to $j$, and the same is
true of $(ia|jb)$. The PNO space of a pair spans the virtual functions of its
own two occupied orbitals, and nothing else. Projecting an amplitude wholesale
into a third pair --- which is what a direct transcription of the canonical
expressions into the local basis amounts to --- therefore truncates the summed
index in a space that has no reason to span it. The resulting error is not a
PNO truncation error: it does not decrease when the PNO threshold is
tightened. Section~\ref{sec:routing} states the rule that avoids it,
classifies every term of the ADC(3) matrix by it, and gives the resulting
evaluation scheme.

The remainder of the paper is organized as follows.
Section~\ref{sec:adc} collects the working equations of IP-ADC($n$) for a
closed-shell reference in the spin-adapted form used here.
Section~\ref{sec:local} constructs the local orbital spaces: localized
occupied orbitals, projected atomic orbitals, differential-overlap domains,
pair natural orbitals and the overlaps that connect them.
Section~\ref{sec:routing} develops the routing rule and applies it term by
term. Section~\ref{sec:amplitudes} treats the first- and second-order
amplitudes, whose local equations are coupled rather than diagonal.
Section~\ref{sec:screening} describes the three-step pair prescreening and the
weak-pair correction, and Section~\ref{sec:fitting} the local density fitting
and the criterion that sets its domains. Section~\ref{sec:solving} gives the
metric and the eigensolver, and Section~\ref{sec:correction} a composite
correction for the truncation of the doubles space.
Section~\ref{sec:implementation} describes the implementation.

% ===========================================================================
\section{IP-ADC for a closed-shell reference}
\label{sec:adc}
% ===========================================================================

We use $i,j,k,l,m,n$ for occupied and $a,b,c,d,e,f$ for virtual spatial
orbitals, and $\pao{\mu},\pao{\nu}$ for projected atomic orbitals. Two-electron
integrals are in chemists' notation, $(ia|jb)=\int
\phi_i(1)\phi_a(1)r_{12}^{-1}\phi_j(2)\phi_b(2)\,d1\,d2$. The Fock matrix in
the occupied block is $F_{ij}$; it is not diagonal once the occupied orbitals
are localized, and this is responsible for much of what follows.

The non-Dyson IP-ADC secular matrix is expressed in the basis of one-hole
(1h) configurations $\ket{i}$ and two-hole--one-particle (2h1p)
configurations $\ket{jka}$,
%
\begin{equation}
  \mathbf{M} =
  \begin{pmatrix} \Mss & \Msd \\[2pt] \Mds & \Mdd \end{pmatrix},
  \qquad
  \mathbf{M}\,\mathbf{X} = -\,\mathbf{X}\,\bm{\Omega},
  \label{eq:secular}
\end{equation}
%
with ionization potentials $\Omega$. In the spin-adapted closed-shell form
the doubles configurations are not orthonormal in the sense required for a
symmetric problem; the appropriate metric is
%
\begin{equation}
  N_{jka,\,lmb} =
    \left(2\,\delta_{jl}\delta_{km} - \delta_{jm}\delta_{kl}\right)\delta_{ab},
  \label{eq:metric}
\end{equation}
%
and the two coupling blocks are related by $\Msd = \mathbf{N}\,\Mds$, so that
only one of them need be constructed. We return to the consequences in
Section~\ref{sec:solving}.

\subsection{The blocks order by order}

Through first order the singles--singles block is the occupied Fock matrix,
%
\begin{equation}
  M^{\text{ss},(0)+(1)}_{ij} = F_{ij},
\end{equation}
%
which in a localized basis is not diagonal. The second-order contribution is
the static part of the self-energy,
%
\begin{align}
  M^{\text{ss},(2)}_{ij} =
    &\sum_{lde}\left(\varepsilon_d + \varepsilon_e - \varepsilon_l\right)
       t^{il}_{de}\,\tilde{t}^{jl}_{de} \nonumber\\
    &+ \left[\,\sum_{lde} K^{il}_{de}\tilde{t}^{jl}_{de}
       - \tfrac12 \varepsilon_i \sum_{lde} t^{il}_{de}\tilde{t}^{jl}_{de}
       + (i \leftrightarrow j)\right],
  \label{eq:mss2}
\end{align}
%
where $K^{ij}_{ab}=(ia|jb)$, $t^{ij}_{ab}$ is the first-order (MP1) amplitude
and $\tilde{t}^{ij}_{ab} = 2t^{ij}_{ab}-t^{ij}_{ba}$. In a localized basis the
orbital energies $\varepsilon_i$, $\varepsilon_l$ appearing in
Eq.~(\ref{eq:mss2}) are replaced by contractions with the occupied Fock
matrix,
%
\begin{equation}
  \varepsilon_i\, t^{il}_{de} \;\longrightarrow\; \sum_k F_{ik}\, t^{kl}_{de},
  \label{eq:eps_sub}
\end{equation}
%
which couples each pair to its neighbours; $\varepsilon_d+\varepsilon_e$
remains diagonal because the virtual functions of each pair are
semicanonicalized within that pair (Section~\ref{sec:local}).

The third-order contribution $M^{\text{ss},(3)}_{ij}$ consists of fourteen
distinct contractions. Written compactly, and with $\hat{t}$ denoting the
second-order amplitude and $t^{(1)\text{v}}$ the particle-ladder contraction
of the first-order amplitude,
%
\begin{align}
  M^{\text{ss},(3)}_{ij} =\;
  &2\sum_{ld} \hat{t}^{(1)}_{ld}\left[2(l\pao{d}|ji)-(j\pao{d}|li)\right]
  \tag{T1}\\
  &+2\sum_{lde}\tilde{\hat{t}}^{\,il}_{de}\,(j\pao{d}|l\pao{e})
  \tag{T2}\\
  &+2\sum_{lde} t^{il}_{de}
     \left(\varepsilon_d+\varepsilon_e-\varepsilon_j-\varepsilon_l\right)
     \tilde{\hat{t}}^{\,jl}_{de}
  \tag{T3}\\
  &+2\sum_{lm}\sum_{def}
     \tilde{t}^{lm}_{de}\,\tilde{t}^{lj}_{df}\,(m\pao{e}|i\pao{f})
  \tag{T4}\\
  &+\sum_{lm}\sum_{def}
     \tilde{t}^{il}_{de}\,(l\pao{e}|m\pao{f})\,\tilde{t}^{jm}_{df}
  \tag{T5}\\
  &-2\sum_{ef} \Theta^{mj}_{ef}\,(mi|\pao{f}\pao{e})
  \tag{T6}\\
  &-\sum_{mef} t^{im}_{fe}
     \Big[\textstyle\sum_{ld}\tilde{t}^{jl}_{fd}(ml|\pao{d}\pao{e})\Big]
  \tag{T7}\\
  &-\sum_{mef} t^{mi}_{fe}
     \Big[\textstyle\sum_{ld}\tilde{t}^{lj}_{fd}(ml|\pao{d}\pao{e})\Big]
  \tag{T8}\\
  &+2\sum_{fe} G^{lm}_{fe}\,(ji|\pao{e}\pao{f})
   \;-\;\sum_{fe} G^{lm}_{fe}\,(j\pao{f}|i\pao{e})
  \tag{T9,T10}\\
  &-2\sum_{mn}\Lambda_{mn}(ji|nm)
   \;+\;\sum_{mn}\Lambda_{mn}(jm|ni)
  \tag{T11,T12}\\
  &+2\sum_{lnm}\Big[\textstyle\sum_{de}t^{in}_{de}\tilde{t}^{lm}_{de}\Big](jl|nm)
  \tag{T13}\\
  &+\sum_{lde}\tilde{t}^{il}_{de}\,t^{(1)\text{v},jl}_{de},
  \tag{T14}
\end{align}
%
with the intermediates
%
\begin{align}
  \Theta^{mj}_{ef} &= \sum_{ld}\left[\,t^{lm}_{de}\tilde{t}^{lj}_{df}
                     + t^{lm}_{ed}\tilde{t}^{jl}_{df}\right], \\
  G^{lm}_{fe}      &= \sum_{d} t^{lm}_{df}\,\tilde{t}^{lm}_{de}, \\
  \Lambda_{mn}     &= \sum_{lde} t^{ln}_{de}\,\tilde{t}^{lm}_{de},\\
  t^{(1)\text{v},ij}_{cd} &= \sum_{ab}(ca|db)\,t^{ij}_{ab}.
\end{align}
%
The occupied indices $l,m,n$ are summed and, together with $i$ and $j$, label
different pairs; it is this structure that Section~\ref{sec:routing} has to
deal with.

The coupling blocks are, through first order,
%
\begin{equation}
  M^{\text{ds},(1)}_{jka,\,i} = (j\pao{a}|ki),
  \qquad
  M^{\text{sd},(1)}_{i,\,jka} = 2(j\pao{a}|ki) - (k\pao{a}|ji),
  \label{eq:mds1}
\end{equation}
%
and acquire at third order
%
\begin{align}
  M^{\text{ds},(3)}_{jka,\,i} =\;
    &\sum_{bc} (i\pao{c}|\pao{a}\pao{b})\, t^{kj}_{cb}
    \tag{D}\\
    &+\sum_{lb} \tilde{t}^{lj}_{ba}\,(l\pao{b}|ik)
    \tag{A}\\
    &-\sum_{lb} t^{lj}_{ba}\,(i\pao{b}|lk)
    \tag{B}\\
    &-\sum_{lb} t^{kl}_{ba}\,(i\pao{b}|lj).
    \tag{C}
\end{align}
%
The doubles--doubles block is diagonal at zeroth order,
%
\begin{equation}
  M^{\text{dd},(0)}_{jka,\,lmb} =
    \left(\varepsilon_j+\varepsilon_k-\varepsilon_a\right)
    \delta_{jl}\delta_{km}\delta_{ab},
  \label{eq:dd0}
\end{equation}
%
--- again with $\varepsilon_j+\varepsilon_k$ replaced by contractions with
$F$ in a localized basis --- and acquires at first order, which is what
distinguishes ADC(3) from ADC(2) in this block,
%
\begin{align}
  \left(\mathbf{M}^{\text{dd},(1)}\mathbf{r}\right)_{jka} =\;
    &-\tfrac12\sum_{il}(ki|jl)\,r_{lia}
      -\tfrac12\sum_{il}(kl|ji)\,r_{ila} \nonumber\\
    &-\sum_{lb}(j\pao{a}|l\pao{b})
       \left[2r_{lkb}-r_{klb}\right] \nonumber\\
    &+\sum_{lb}(jl|\pao{b}\pao{a})\,r_{lkb}
      +\sum_{lb}(kl|\pao{b}\pao{a})\,r_{jlb}.
  \label{eq:dd1}
\end{align}
%
The two contributions on the first line are the same contraction written twice
over complementary index ranges and are evaluated together.

% ===========================================================================
\section{Local orbital spaces}
\label{sec:local}
% ===========================================================================

\subsection{Occupied and virtual bases}

The occupied orbitals are localized by the Pipek--Mezey
criterion.\cite{PipekMezey1989} Frozen core orbitals are excluded from the
localization and from the correlation treatment. The virtual space is spanned
by projected atomic orbitals,
%
\begin{equation}
  \ket{\pao{\mu}} = \Big(1 - \sum_{i}^{\text{all occ}}\ket{i}\bra{i}\Big)\ket{\mu},
  \label{eq:pao}
\end{equation}
%
where the projector runs over \emph{all} occupied orbitals including the
frozen core; projecting only the correlated occupied space would leave core
density in the virtual space and inflate every domain. PAOs are non-orthogonal
and linearly dependent, with overlap $S^{\pao{\mu}\pao{\nu}}$.

\subsection{Domains}

Locality is measured by the differential overlap
integral\cite{Pinski2015}
%
\begin{equation}
  \mathrm{DOI}(i,\pao{\mu}) =
    \left[\int \left|\phi_i(\mathbf{r})\right|^2
                \left|\phi_{\pao{\mu}}(\mathbf{r})\right|^2
          d\mathbf{r}\right]^{1/2},
  \label{eq:doi}
\end{equation}
%
evaluated on a numerical grid. An atom $A$ belongs to the domain of $i$ if any
PAO centred on $A$ satisfies $\mathrm{DOI}(i,\pao{\mu}) > \TCut{DO}$, and the
domain $[i]$ consists of all PAOs on the atoms so selected. Domains are
atom-granular by construction: a domain that cut through an atom's basis set
would not be invariant to rotations of that basis. The domain of a pair is the
union $[ij] = [i]\cup[j]$.

Within a pair domain the PAOs are made orthonormal and semicanonical in two
steps. The domain overlap is diagonalized and eigenvectors with eigenvalues
below $\TCut{PAO}$ are discarded, removing the linear dependence; the Fock
matrix is then diagonalized in the surviving space. We call the result the
\emph{non-redundant semicanonical domain basis} of the pair, of dimension
$n^{ij}_{\text{nr}}$, and denote its transformation matrix $\mathbf{d}^{ij}_{\text{nr}}$.
It is the largest space the pair can represent, and it will be needed in
Section~\ref{sec:amplitudes}.

\subsection{Pair natural orbitals}

In that basis the semicanonical first-order amplitude is
%
\begin{equation}
  \bar{t}^{ij}_{ab} = \frac{-K^{ij}_{ab}}
    {\varepsilon_a + \varepsilon_b - F_{ii} - F_{jj}},
  \label{eq:tsc}
\end{equation}
%
and the pair density is
%
\begin{equation}
  \mathbf{D}^{ij} = \frac{2}{1+\delta_{ij}}
    \left[ \tilde{\bar{\mathbf{t}}}^{\,\dagger}\bar{\mathbf{t}}
         + \tilde{\bar{\mathbf{t}}}\,\bar{\mathbf{t}}^{\dagger}\right].
  \label{eq:paird}
\end{equation}
%
Its eigenvectors with occupation numbers above $\TCut{PNO}$ are retained as
the pair natural orbitals of $(i,j)$; the Fock matrix is diagonalized once more
within the retained space so that $\varepsilon_a$ remains a genuine orbital
energy there. Writing $\mathbf{v}^{ij}$ for the retained eigenvectors, the PNO
coefficients in the domain are
%
\begin{equation}
  \mathbf{d}^{ij} = \mathbf{d}^{ij}_{\text{nr}}\,\mathbf{v}^{ij},
  \label{eq:dpno}
\end{equation}
%
a factorization we use repeatedly: $\mathbf{d}^{ij}_{\text{nr}}$ spans the
whole domain, $\mathbf{d}^{ij}$ only the part of it the pair correlates into.
The unordered pair $\{i,j\}$ carries one PNO space, shared by the two ordered
labels $(i,j)$ and $(j,i)$, since the density~(\ref{eq:paird}) is symmetric
under the exchange even though the amplitude is not.

\subsection{Cross-pair overlaps}

Quantities belonging to different pairs are related by the overlap of their
PNO spaces,
%
\begin{equation}
  S^{PQ}_{ab} = \big\langle a_P \big| b_Q \big\rangle
    = \big[(\mathbf{d}^{P})^{\dagger}\,\mathbf{S}^{\text{PAO}}\,
           \mathbf{d}^{Q}\big]_{ab},
  \label{eq:spq}
\end{equation}
%
where $\mathbf{S}^{\text{PAO}}$ is restricted to the union of the two domains.
A matrix $\mathbf{X}$ living in pair $Q$ is expressed in pair $P$ as
$\mathbf{S}^{PQ}\mathbf{X}(\mathbf{S}^{PQ})^{\dagger}$. This transformation is
exact if and only if $\mathbf{X}$ has no support outside the space
$\mathbf{d}^{P}$ spans --- a condition that is the subject of the next section.

% ===========================================================================
\section{Evaluating the ADC matrix in the local basis}
\label{sec:routing}
% ===========================================================================

\subsection{The rule}

Consider a contraction of the generic form
%
\begin{equation}
  X = \sum_{c} g_{\ldots c \ldots}\; A^{Q}_{\ldots c \ldots},
  \label{eq:generic}
\end{equation}
%
in which a virtual index $c$ is summed between a two-electron integral $g$ and
an amplitude $A$ belonging to pair $Q$, and the result $X$ is required in the
space of some other pair $P$. Two evaluations suggest themselves. One may
transform $A$ into $P$ first and contract there, which requires the integral
only in $P$'s space; or one may contract in $Q$, where $A$ is exact, and
transform the result.

The two are not equivalent, and the choice is not a matter of taste. Every
virtual index in Section~\ref{sec:adc} is associated with a definite occupied
label --- $a$ with $i$ and $b$ with $j$ in both $t^{ij}_{ab}$ and $(ia|jb)$ ---
and the PNO space of a pair is constructed, through
Eqs.~(\ref{eq:tsc})--(\ref{eq:paird}), to span the virtual functions its own
two occupied orbitals correlate into. It has no claim on the virtual functions
of a third orbital. We therefore adopt the following rule.

\medskip
\noindent\textbf{Rule.} \emph{A virtual index summed between a two-electron
integral and an amplitude must be carried in a pair space that owns it, i.e.\
one whose occupied labels include the occupied label the index belongs to.}
\medskip

Two classes of contraction are exempt, and it is useful to name them.

\emph{(i) Amplitude--amplitude contractions.} If the summed index connects two
amplitudes with no integral between them, both factors decay with the index,
and the projection $\mathbf{S}^{PQ}\mathbf{X}(\mathbf{S}^{PQ})^{\dagger}$ into
either pair is accurate to the PNO truncation of the amplitudes themselves.
Terms T11--T13 and the $\Lambda_{mn}$, $G^{lm}$ intermediates are of this kind.

\emph{(ii) Coulomb-type integrals.} If the two virtual indices of an integral
sit on the same electron, as in $(xy|\pao{a}\pao{b})$, the integral does not
assign them to separate occupied labels and both must in any case share a
space. The choice is then determined by the surrounding amplitudes rather than
by the integral.

Everything else falls under the rule, and violating it has a characteristic
signature: because the neglected part of the summand is not the tail of a
converging PNO expansion but a piece of a different pair's space, the error
does \emph{not} decrease as $\TCut{PNO}$ is tightened. This is what makes the
mistake easy to miss --- the calculation looks converged.

\subsection{Term-by-term classification}

We now apply the rule to each block. Throughout, $P$ denotes the target pair
and $\rightarrow$ a transformation by $\mathbf{S}^{PQ}$.

\paragraph{Terms native to one pair.}
T2 contracts $\tilde{\hat{t}}^{\,il}$ with $(j\pao{d}|l\pao{e})$, which is the
$ovov$ block of the pair $(i,l)$ itself; T9 and T10 contract two amplitudes of
the same pair $(l,m)$ with integrals of that pair; the particle ladder in
$t^{(1)\text{v}}$ never leaves the pair it belongs to. These are evaluated in
place, with no transformation at all.

\paragraph{Terms of class (i).}
T3 contracts $t^{il}$ with a function of $\tilde{\hat{t}}^{\,jl}$; both carry
the index $l$, and the transformation $(j,l)\rightarrow(i,l)$ is exact in the
sense of (i). T11--T13 likewise. The occupied-Fock substitutions of
Eq.~(\ref{eq:eps_sub}) are of the same character.

\paragraph{T4 and T6: intermediates whose indices live in different pairs.}
In T4 the index $d$ is summed between $\tilde{t}^{lm}$ and $\tilde{t}^{lj}$
and belongs to $l$; $e$ belongs to $m$ and $f$ to $j$, while the integral
$(m\pao{e}|i\pao{f})$ belongs to the pair $(m,i)$. Building the intermediate
in $(m,j)$ --- the choice suggested by its occupied labels, and the one that
would allow the $l$ sum to be accumulated once and reused for every $i$ ---
places both $d$ and $f$ in a pair that owns neither. Instead we form
%
\begin{equation}
  \mathcal{G}^{(l,m),(l,j)}_{ef} = \sum_{d}
    \tilde{t}^{lm}_{de}\left[\mathbf{S}^{(lm)(lj)}\tilde{\mathbf{t}}^{lj}\right]_{df},
  \label{eq:t4}
\end{equation}
%
whose two indices lie in \emph{different} pairs, each of which owns its index,
and only then map it into the integral's pair,
$\mathcal{G} \rightarrow \mathbf{S}^{(mi)(lm)}\,\mathcal{G}\,
(\mathbf{S}^{(mi)(lj)})^{\dagger}$. The price is that the $l$ summation can no
longer be pre-accumulated and the loop is one factor of $n_{\text{occ}}$
deeper; the benefit is that the term converges. T6 is treated identically,
with the second half of $\Theta^{mj}$ carrying $e$ in $(l,m)$ and $f$ in
$(j,l)$.

\paragraph{T5: a contested index contracted directly.}
Here $e$ belongs to $l$ and $f$ to $m$, so both may be brought to the
integral's pair $(l,m)$, which owns them. The index $d$, however, belongs to
$i$ in one amplitude and to $j$ in the other, and $(l,m)$ owns neither. It is
therefore never represented there: with
%
\begin{align}
  \mathcal{A}^{i}_{df} &= \left[\tilde{\mathbf{t}}^{il}\,
     \mathbf{S}^{(il)(lm)}\,\mathbf{K}^{lm}\right]_{df}, \\
  \mathcal{B}^{j}_{df} &= \left[\tilde{\mathbf{t}}^{jm}\,
     \mathbf{S}^{(jm)(lm)}\right]_{df},
\end{align}
%
in which only the index the integral owns has been mapped, the term is
%
\begin{equation}
  \text{T5} = \sum_{lm}
    \Big\langle \mathcal{A}^{i}(\mathcal{B}^{j})^{\dagger},\;
                \mathbf{S}^{(il)(jm)}\Big\rangle,
  \label{eq:t5}
\end{equation}
%
the contested index being contracted directly between the two amplitudes'
own spaces.

\paragraph{T7 and T8: choosing the integral's home.}
The integral $(ml|\pao{d}\pao{e})$ is of class (ii), but the amplitudes around
it still tie $d$ to $l$ and $e$ to $m$. The pair that owns both is $(l,m)$,
and that is where it is read: it is that pair's own $oovv$ block, indexed by
its own two occupied labels, so it always exists and requires no screening of
occupied lists. From there $d$ moves to $(j,l)$ to meet $\tilde{t}^{jl}$ and
$e$ to $(i,m)$ to meet $t^{im}$, each sharing an occupied label with its
destination. The remaining index $f$ is genuinely contested --- it belongs to
$j$ in the amplitude and to $i$ in the target --- and is mapped once; that
mapping is unavoidable and is what makes the term small when $i$ and $j$ are
far apart.

\paragraph{$\mathbf{M}^{\text{ds},(3)}$.}
Term D keeps every index in the target pair $(j,k)$, which owns them all. In
terms A, B and C the summed index $b$ belongs to $l$ (A, B) or to $k$ (C), and
the target has no claim on it: the contraction over $b$ is therefore performed
in the \emph{amplitude's} pair, reading the $ooov$ integral out of that pair's
block, and only the surviving index $a$ --- which both pairs own, since they
share an occupied label --- is carried across with $\mathbf{S}^{PQ}$.

\paragraph{The first-order doubles block.}
In Eq.~(\ref{eq:dd1}) the two $oooo$ contractions involve no integral over
virtual indices, so transformation into the target pair is exact in the sense
of (i). The remaining three do: in $(j\pao{a}|l\pao{b})$ the index $b$ belongs
to $l$ and is summed against $r_{lkb}$, and the target $(j,k)$ has no claim on
it. Each is therefore evaluated in the pair that owns $b$ --- $(j,l)$ for the
exchange and the first Coulomb line, $(k,l)$ for the second --- and mapped
back afterwards. Two transformations replace one, but both are between pairs
sharing an occupied index with the index being moved.

\paragraph{Screening.} The routing determines the screening rather than the
other way round. In T4, for example, $l$ contributes only if $(l,m)$ and
$(l,j)$ are both strong pairs, and $i$ only if $(m,i)$ is one; no pair
$(i,j)$ is ever required, so $M^{\text{ss}}_{ij}$ survives for $i,j$ that do
not themselves form a strong pair, as it must.

% ===========================================================================
\section{Amplitudes in the local basis}
\label{sec:amplitudes}
% ===========================================================================

\subsection{First-order amplitudes}

Because $F_{ij}$ is not diagonal, the first-order amplitudes satisfy a coupled
system rather than a division,
%
\begin{align}
  R^{ij}_{ab} = K^{ij}_{ab}
    &- \left(F_{ii}+F_{jj}-\varepsilon_a-\varepsilon_b\right) t^{ij}_{ab}
      \nonumber\\
    &- \sum_{k \ne i} F_{ik}
        \left[\mathbf{S}\,\mathbf{t}^{kj}\mathbf{S}^{\dagger}\right]_{ab}
      - \sum_{k \ne j} F_{jk}
        \left[\mathbf{S}\,\mathbf{t}^{ik}\mathbf{S}^{\dagger}\right]_{ab},
  \label{eq:lmp2}
\end{align}
%
solved iteratively with the semicanonical denominator as preconditioner. The
coupling is of class (i) and the transformations are exact.

\subsection{Second-order amplitudes}

The second-order amplitude obeys the same coupled system with a different
right-hand side,
%
\begin{align}
  \mathcal{D}^{ij}_{ab} =\;
    &\sum_{cd}(ca|db)\,t^{ij}_{cd}
     + \sum_{kl}(ki|lj)\,t^{kl}_{ab} \nonumber\\
    &+ 2\sum_{kc}(kc|bj)\,t^{ki}_{ca}
     - \sum_{kc}(kc|bj)\,t^{ik}_{ca} \nonumber\\
    &- \sum_{kc}(kj|bc)\,t^{ik}_{ac}
     - \sum_{kc}(ki|bc)\,t^{kj}_{ac}
     - \sum_{kc}(kj|ac)\,t^{ik}_{cb} \nonumber\\
    &+ 2\sum_{kc}(kc|ai)\,t^{kj}_{cb}
     - \sum_{kc}(kc|ai)\,t^{jk}_{cb}
     - \sum_{kc}(ki|ac)\,t^{kj}_{cb}.
  \label{eq:t2}
\end{align}
%
Solving Eq.~(\ref{eq:t2}) by division, as the canonical expression permits,
discards the inter-pair Fock coupling and is a measurable error; the coupled
system is solved instead.

The first two terms are safe --- the particle ladder never leaves the pair,
and the $oooo$ term sums no virtual against an integral --- and are evaluated
in the pair's PNO space. The remaining eight are the ring terms, and each of
them sums a virtual $c$ between an integral and an amplitude $t^{Q}$ that
belongs to a \emph{neighbouring} pair $Q$, which is exactly the situation
Section~\ref{sec:routing} legislates for. It is worth being precise about
which pair owns $c$, because the obvious candidates are not the right one.

The sharpest answer is $Q$ itself. The amplitude is identically zero outside
$\mathrm{PNO}(Q)$, so
%
\begin{equation}
  \sum_{c \in \mathrm{PNO}(Q)} g_{\ldots c \ldots}\, t^{Q}_{\ldots c \ldots}
  \;=\; \sum_{c} g_{\ldots c \ldots}\, t^{Q}_{\ldots c \ldots}
  \label{eq:ampspace}
\end{equation}
%
is not an approximation to the full-space sum but that sum itself, to the
accuracy of the amplitude, which is the truncation the method is built on in
any case. The two indices that survive are projected into the target pair's
PNO space, where the driver lives; only the index that has to travel is
transformed, and the amplitude is never projected. What this costs is a mixed
integral, one virtual in $Q$ and one in the target pair. It costs no storage:
$(k\pao{c}_{Q}|\pao{b}_{P}x)$ is the two pairs' own three-index blocks
contracted over a common fitting domain, and
$(kx|\pao{b}_{P}\pao{c}_{Q})$ is formed by contracting the auxiliary index
against the occupied pair first, so that the mixed $(\mu\nu)$ block appears
once and the two virtual transforms are matrix multiplications.

Two spaces that look adequate are not. Carrying $c$ in the target pair's own
non-redundant domain basis $\mathbf{d}^{ij}_{\text{nr}}$ of
Eq.~(\ref{eq:dpno}) is generous --- it spans the whole domain of $(i,j)$ ---
but $[i]\cup[j]$ need not contain $[k]$, and it is near $k$ that the integral
is largest. Carrying it in the orbital domain $[k]$ of the summed occupied
fixes that and breaks the other side, since $[k]$ need not contain
$\mathrm{PNO}(Q)$. Table~\ref{tab:ringspace} measures the three on a chain of
four water molecules $3.2\,\text{\AA}$ apart, as the \emph{domain} error of
the MP3 energy: the same code at $\text{TCutDO}=10^{-2}$ minus the same code
with the domains switched off, so that the PNO truncation, which all three
share, divides out. The two domain answers bracket the exact one from opposite
sides --- the signature of an amplitude being projected rather than an
integral being screened --- while Eq.~(\ref{eq:ampspace}) leaves no measurable
domain error at all. It is also the cheapest of the three, a PNO pair being a
handful of functions where a PAO domain is tens.

\begin{table}
  \caption{Domain error of the MP3 energy from the eight ring terms of
    Eq.~(\ref{eq:t2}), in $\mu E_h$, for three choices of the space carrying
    the summed virtual $c$. $(\text{H}_2\text{O})_4$ chain,
    $3.2\,\text{\AA}$ spacing, cc-pVDZ.}
  \label{tab:ringspace}
  \begin{tabular}{lrrr}
    \hline\hline
    \TCut{PNO} & $\mathrm{PNO}(Q)$ & $[i]\cup[j]$ & $[k]$ \\
    \hline
    $10^{-5}$ & $-0.15$ & $-4.88$  & $+5.51$  \\
    $10^{-6}$ & $-0.43$ & $-50.66$ & $+43.63$ \\
    \hline\hline
  \end{tabular}
\end{table}

A molecule whose domains are complete cannot see any of this: every candidate
space is then the whole virtual space and all three are exact. Water is such a
molecule, which is why the measurement above required a chain.

\subsection{Second-order singles}

Term T1 of $M^{\text{ss},(3)}$ requires the second-order singles amplitude,
whose driver is
%
\begin{equation}
  \mathcal{D}^{(1)}_{ia} = \sum_{kcd}\tilde{t}^{ik}_{cd}\,(k\pao{d}|\pao{a}\pao{c})
    - \sum_{ick}(i\pao{c}|kj)\,\tilde{t}^{ik}_{cd}\Big|_{jd},
  \label{eq:t1drv}
\end{equation}
%
and which is represented in the PNO space of the diagonal pair $(i,i)$. As
with Eq.~(\ref{eq:lmp2}) the equation couples different $i$ through $F$ and is
solved as a linear system.

% ===========================================================================
\section{Pair prescreening and the weak-pair correction}
\label{sec:screening}
% ===========================================================================

A pair that is classified strong carries a PNO space, an amplitude, four
integral blocks and a slice of the ADC eigenvector; a pair that is not costs
nothing. The classification is made in three steps.

\paragraph{Step 1: a cheap estimate over every pair.}
Two channels are used, because no single inexpensive quantity sees both
sources of a pair energy.

The \emph{overlap channel} is the one that decides on ordinary molecules. When
two localized orbitals share space the pair energy is driven by their exchange
integral, which involves the product density $\phi_i(\mathbf{r})
\phi_j(\mathbf{r})$; the natural scaling is therefore quadratic in the
differential overlap,
%
\begin{equation}
  e_{ij} \;\approx\; -\,c_{\text{DOI}}\;\mathrm{DOI}(i,j)^2 .
  \label{eq:doiest}
\end{equation}
%
Measured over every occupied pair of a four-water chain at 3.2~\AA{} and of a
C$_8$H$_{16}$ chain, both in cc-pVDZ, the ratio
$|e_{ij}|/\mathrm{DOI}(i,j)^2$ lies between $0.3$ and $10^{2}$ across all
pairs worth more than $10^{-6}$~$E_h$ --- two orders of magnitude of scatter
for pair energies spanning five. The smallest prefactor that would retain
every pair worth more than $\TCut{Pairs}$, at $\TCut{Pre}=\TCut{Pairs}/100$,
is $0.38$ on those systems; we use $c_{\text{DOI}}=25$, a safety factor of
roughly sixty, and the resulting estimate of the discarded correlation energy
agrees with the true value to within about ten percent.

The \emph{dispersion channel} covers what Eq.~(\ref{eq:doiest}) cannot see.
Two orbitals that do not overlap at all still have an $R^{-6}$ pair energy,
and the differential overlap is exponentially blind to it. There the exchange
integral is dominated by the dipole--dipole interaction of the two transition
densities, giving the semicanonical estimate
%
\begin{equation}
  e^{\text{dip}}_{ij} = -\frac{4}{R_{ij}^{6}}
    \sum_{a\in[i]}\sum_{b\in[j]}
    \frac{\left[\bm{\mu}_{ia}\!\cdot\!\bm{\mu}_{jb}
      - 3(\bm{\mu}_{ia}\!\cdot\!\hat{\mathbf{R}})
         (\bm{\mu}_{jb}\!\cdot\!\hat{\mathbf{R}})\right]^{2}}
    {\varepsilon_a+\varepsilon_b-F_{ii}-F_{jj}},
  \label{eq:dipole}
\end{equation}
%
with $\bm{\mu}_{ia}=\bra{i}\mathbf{r}\ket{\pao{a}}$ and $\mathbf{R}_{ij}$ the
vector between the centroids of $i$ and $j$; $a$ and $b$ run over each
orbital's own domain, orthogonalized and semicanonicalized, so that
$\varepsilon_a$ is a genuine orbital energy of that domain. Both the
transition dipoles and the centroid difference are origin-independent, the
former because the PAO projector~(\ref{eq:pao}) makes PAOs orthogonal to all
occupied orbitals. Equation~(\ref{eq:dipole}) is formed only for pairs whose
differential overlap lies below $\TCut{DO}_{ij}$, since a multipole expansion
between overlapping distributions is meaningless.

A pair is retained if \emph{either} channel exceeds $\TCut{Pre}$ in magnitude;
neither is a veto. In practice the overlap channel does nearly all of the
work --- on the systems above the dispersion channel alters one decision in
several hundred --- but it is inexpensive and the situation it covers, a pair
of well-separated polarizable fragments, is a real one.

The estimate is evaluated for every unordered pair rather than for a
pre-filtered subset. It is $O(1)$ per pair, so there is nothing to be gained
by filtering ahead of it, and a topological filter --- retaining only pairs
whose domains share an atom, say --- would discard pairs without an energy
attached to them, leaving a hole in the accounting of
Eq.~(\ref{eq:eweak}) below.

\paragraph{Step 2: semicanonical pair energies.}
For every survivor the semicanonical local MP2 pair energy is computed in the
pair's own domain,
%
\begin{equation}
  e_{ij} = (2-\delta_{ij})\sum_{ab} K^{ij}_{ab}
    \left(2\bar{t}^{ij}_{ab}-\bar{t}^{ij}_{ba}\right),
  \label{eq:pairE}
\end{equation}
%
with $\bar{t}$ from Eq.~(\ref{eq:tsc}). This is the quantity $\TCut{Pairs}$ is
applied to. It costs one domain orthogonalization, one local fit and one
division per candidate --- the same work as the first steps of the PNO
construction, without the PNOs --- and it is a far better criterion than
Eq.~(\ref{eq:dipole}), which is blind to overlap-driven pair energies by
construction.

\paragraph{Step 3: the weak-pair correction.}
The pair energies of the discarded pairs are summed,
%
\begin{equation}
  E_{\text{weak}} = \sum_{ij\,\in\,\text{weak}} e_{ij}
    + \sum_{ij\,\in\,\text{step 1}} e^{\text{est}}_{ij},
  \label{eq:eweak}
\end{equation}
%
using the best estimate available for each: the semicanonical pair energy for
pairs step 2 rejected, and whichever of the two step-1 channels was larger in
magnitude for those rejected earlier. $E_{\text{weak}}$ is a real approximation to the
correlation energy discarded and is reported alongside the local MP2 energy;
their sum estimates the unscreened result.

% ===========================================================================
\section{Local density fitting}
\label{sec:fitting}
% ===========================================================================

All two-electron integrals are obtained from a density-fitting
approximation,\cite{Whitten1973,Vahtras1993}
%
\begin{equation}
  (pq|rs) \approx \sum_{KL} (pq|K)\,[\mathbf{J}^{-1}]_{KL}\,(L|rs),
\end{equation}
%
in which the auxiliary index is restricted to a fitting domain that depends
on the pair. The functions a pair's fit must represent are the products
$\phi_i(\mathbf{r})\phi_{\pao{\mu}}(\mathbf{r})$, whose extent is governed by
the \emph{occupied} orbital; the PAO domain answers a different question and is
not the right criterion. We therefore select fitting atoms from the orbital's
own atomic populations, in the spirit of Mata and
Werner:\cite{MataWerner2006} for each localized occupied orbital the atoms are
ordered by decreasing L\"owdin population of $|\phi_i|^2$ and taken in that
order until the population left behind falls below $\TCut{MKN}$; the fitting
domain of a pair is the union of its two orbitals' atom sets, and the
auxiliary functions on them. A L\"owdin rather than a Mulliken partition is
used because it is positive, so that ``the population left behind'' is a
quantity a threshold can meaningfully be applied to. Smaller $\TCut{MKN}$
retains more atoms, in the same direction as every other threshold in the
method.

A local fitting domain is a subset of an auxiliary basis that was never
designed to be complete on its own, so the restricted Coulomb metric is
routinely ill-conditioned. Its inverse is therefore formed as a pseudo-inverse
from the eigendecomposition, with directions below a relative cutoff
discarded; a Cholesky or LU factorization does not fail on such a matrix but
returns a fit that is wrong in the tails.

% ===========================================================================
\section{The metric and the eigensolver}
\label{sec:solving}
% ===========================================================================

The metric~(\ref{eq:metric}) is block diagonal over unordered pairs in the PNO
basis, each block acting on the pair of ordered labels $(j,k)$ and $(k,j)$ as
$2\mathbf{1}-\mathbf{P}$ with $\mathbf{P}$ the exchange of the two. Its
eigenvalues are therefore $1$ on the symmetric combination and $3$ on the
antisymmetric one, and $\mathbf{N}^{\pm1/2}$ is available in closed form. The
similarity transformation
%
\begin{equation}
  \mathbf{N}^{-1/2}\,\mathbf{M}\,\mathbf{N}^{-1/2}\,\mathbf{y}
    = -\Omega\,\mathbf{y},
  \qquad \mathbf{y} = \mathbf{N}^{1/2}\mathbf{x},
\end{equation}
%
renders the problem symmetric at no cost, and a symmetric Davidson
procedure\cite{Davidson1975} is used, with the diagonal of $\Mdd$ as
preconditioner --- which remains diagonal precisely because the PNOs of each
pair are semicanonicalized.

One caveat is worth recording. The third-order static self-energy as
conventionally written is not symmetric, although the exact self-energy is
Hermitian; the asymmetry is an omission in the working equations rather than
physics. It is small --- it moves the roots by well under 1~meV --- but it
breaks the self-adjointness the symmetric Davidson relies on, so that the
residual cannot fall below the asymmetry and the iteration stalls short of its
tolerance. $\mathbf{M}^{\text{ss},(3)}$ is therefore symmetrized before diagonalization.

% ===========================================================================
\section{A composite correction for the doubles space}
\label{sec:correction}
% ===========================================================================

The dominant approximation of the method is the truncation of the 2h1p space
to the PNOs of each pair. At strict second order this error can be measured
rather than estimated. The 2h1p--2h1p block of IP-ADC(2) is
Eq.~(\ref{eq:dd0}) alone, so L\"owdin partitioning of
Eq.~(\ref{eq:secular}) is an identity, not an approximation:\cite{Mester2023}
%
\begin{equation}
  \left[\Mss + \mathbf{F}(\omega)\right]\mathbf{c} = \omega\,\mathbf{c},
  \qquad
  \mathbf{F}(\omega) = \Msd\left(\omega-\Mdd\right)^{-1}\Mds ,
  \label{eq:fold}
\end{equation}
%
with $\Omega=-\omega$ and no doubles vector ever formed. Canonically
$(\omega-\Mdd)^{-1}$ is a scalar division; in the PNO basis only the
$\varepsilon_a$ part is diagonal and the two occupied indices become Fock
couplings to neighbouring pairs, so Eq.~(\ref{eq:fold}) becomes a shifted
linear system in the doubles space, which is solved rather than approximated
--- dropping the coupling would contaminate the difference we are trying to
measure.

Running the fold canonically, before localization, and again in the PNO basis
gives the second-order truncation error directly. We use the
scaled-opposite-spin variant of the second-order self-energy for this purpose,
in which the eight terms of Eq.~(\ref{eq:mss2}) collapse to two scaled by
$c_{\text{os}}$, the coupling blocks and $\Mdd$ being untouched. The
difference is added to the third-order result as a composite correction. It
removes the doubles-space channel exactly; it cannot touch the third-order
amplitude channel, and whether the composite is an improvement is therefore an
empirical question.

% ===========================================================================
\section{Implementation}
\label{sec:implementation}
% ===========================================================================

The method is implemented in the BAGH program package. It is dispatched on the
method name before any canonical integral transformation is performed, since
it reads nothing from the conventional integral machinery: letting that run
first would carry out precisely the $O(N^5)$ transformation --- including, at
the default settings, the dense $(vv|vv)$ block --- that the local treatment
exists to avoid. Orbital localization,
the projected atomic orbitals, the differential-overlap and dipole matrices
and the three-index integrals are prepared in Python on top of
PySCF;\cite{Sun2018} everything after that --- domains, prescreening, pair
natural orbitals, the integral store, the amplitude equations, the
intermediates and the Davidson procedure --- is carried out in C++ behind a
pybind11 interface, parallelized with OpenMP. No trial vector, sigma vector or
amplitude crosses the language boundary.

Sparsity is represented throughout by sparse maps in the sense of Pinski
\emph{et al.}\cite{Pinski2015} All two-electron integrals are generated once,
after the PNO spaces are complete, directly in the PNO basis of the pair that
owns them; nothing is generated on the fly and no integral is ever stored in a
canonical or PAO basis. For each unordered pair $U$ the store holds
$(x\pao{a}|y\pao{b})$, $(xy|\pao{a}\pao{b})$, $(x\pao{a}|yz)$ and
$(x\pao{a}|\pao{b}\pao{c})$, the occupied labels $x,y,z$ running over a list
comprising the pair's own two labels together with those whose differential
overlap with the pair's domain exceeds a threshold.

Because the largest of these classes can exceed available memory while the
smallest cannot, residency is planned rather than left to the operating
system. Before any integral is formed, the size of each class is computed from
the PNO dimensions and compared against the user's memory limit; classes are
assigned to memory or to a memory-mapped file in order of access frequency,
with the pair data, the Davidson subspace and the $oooo$ block never spilled.
A calculation whose unspillable data exceeds the limit stops with a report
rather than paging.

A second, independent implementation of the same equations is maintained
alongside the first. It is written in dense linear algebra with no sparse
maps, no memory management and no Davidson procedure: every quantity is a
dense array, each contraction is a single tensor contraction, and the full
ADC matrix is formed and diagonalized directly. It shares no code with the
production implementation, and reproduces it to $10^{-3}$~meV in the
ionization potentials and $10^{-12}$~$E_h$ in the second- and third-order
energies at every threshold tested. Its purpose is to check the working
equations rather than their programming, and it was in fact the means by which
two of the routing choices of Section~\ref{sec:routing} were found to be
suboptimal in the production code.

% ===========================================================================
\section*{Author declarations}
The authors have no conflicts to disclose.

\section*{Data availability}
The data that support the findings of this study are available from the
corresponding author upon reasonable request.

\begin{thebibliography}{99}

\bibitem{Schirmer1982} J. Schirmer,
  \emph{Beyond the random-phase approximation: A new approximation scheme for
  the polarization propagator}, Phys. Rev. A \textbf{26}, 2395 (1982).

\bibitem{Schirmer1983} J. Schirmer, L. S. Cederbaum and O. Walter,
  \emph{New approach to the one-particle Green's function for finite Fermi
  systems}, Phys. Rev. A \textbf{28}, 1237 (1983).

\bibitem{Trofimov2005} A. B. Trofimov and J. Schirmer,
  \emph{Molecular ionization energies and ground- and ionic-state properties
  using a non-Dyson electron propagator approach},
  J. Chem. Phys. \textbf{123}, 144115 (2005).

\bibitem{Dreuw2015} A. Dreuw and M. Wormit,
  \emph{The algebraic diagrammatic construction scheme for the polarization
  propagator for the calculation of excited states},
  WIREs Comput. Mol. Sci. \textbf{5}, 82 (2015).

\bibitem{Pulay1983} P. Pulay,
  \emph{Localizability of dynamic electron correlation},
  Chem. Phys. Lett. \textbf{100}, 151 (1983).

\bibitem{Saebo1993} S. Saeb{\o} and P. Pulay,
  \emph{Local treatment of electron correlation},
  Annu. Rev. Phys. Chem. \textbf{44}, 213 (1993).

\bibitem{Neese2009} F. Neese, F. Wennmohs and A. Hansen,
  \emph{Efficient and accurate local approximations to coupled-electron pair
  approaches}, J. Chem. Phys. \textbf{130}, 114108 (2009).

\bibitem{Riplinger2013a} C. Riplinger and F. Neese,
  \emph{An efficient and near linear scaling pair natural orbital based local
  coupled cluster method}, J. Chem. Phys. \textbf{138}, 034106 (2013).

\bibitem{Riplinger2013b} C. Riplinger, B. Sandhoefer, A. Hansen and F. Neese,
  \emph{Natural triple excitations in local coupled cluster calculations with
  pair natural orbitals}, J. Chem. Phys. \textbf{139}, 134101 (2013).

\bibitem{Pinski2015} P. Pinski, C. Riplinger, E. F. Valeev and F. Neese,
  \emph{Sparse maps --- A systematic infrastructure for reduced-scaling
  electronic structure methods. I. An efficient and simple linear scaling
  local MP2 method}, J. Chem. Phys. \textbf{143}, 034108 (2015).

\bibitem{Riplinger2016} C. Riplinger, P. Pinski, U. Becker, E. F. Valeev and
  F. Neese, \emph{Sparse maps --- A systematic infrastructure for
  reduced-scaling electronic structure methods. II. Linear scaling domain
  based pair natural orbital coupled cluster theory},
  J. Chem. Phys. \textbf{144}, 024109 (2016).

\bibitem{Dutta2018} A. K. Dutta, M. Saitow, C. Riplinger, F. Neese and
  R. Izs\'ak, \emph{A near-linear scaling equation of motion coupled cluster
  method for ionized states},
  J. Chem. Phys. \textbf{148}, 244101 (2018).

\bibitem{PipekMezey1989} J. Pipek and P. G. Mezey,
  \emph{A fast intrinsic localization procedure applicable for ab initio and
  semiempirical linear combination of atomic orbital wave functions},
  J. Chem. Phys. \textbf{90}, 4916 (1989).

\bibitem{MataWerner2006} R. A. Mata and H.-J. Werner,
  \emph{Calculation of smooth potential energy surfaces using local electron
  correlation methods}, J. Chem. Phys. \textbf{125}, 184110 (2006).

\bibitem{Whitten1973} J. L. Whitten,
  \emph{Coulombic potential energy integrals and approximations},
  J. Chem. Phys. \textbf{58}, 4496 (1973).

\bibitem{Vahtras1993} O. Vahtras, J. Alml\"of and M. W. Feyereisen,
  \emph{Integral approximations for LCAO-SCF calculations},
  Chem. Phys. Lett. \textbf{213}, 514 (1993).

\bibitem{Davidson1975} E. R. Davidson,
  \emph{The iterative calculation of a few of the lowest eigenvalues and
  corresponding eigenvectors of large real-symmetric matrices},
  J. Comput. Phys. \textbf{17}, 87 (1975).

\bibitem{Mester2023} D. Mester and M. K\'allay,
  \emph{Vertical ionization potentials and electron affinities at the
  double-hybrid density functional level},
  J. Chem. Theory Comput. \textbf{19}, 3982 (2023).

\bibitem{Sun2018} Q. Sun \emph{et al.},
  \emph{PySCF: the Python-based simulations of chemistry framework},
  WIREs Comput. Mol. Sci. \textbf{8}, e1340 (2018).

\end{thebibliography}

\end{document}
